\chapter{Choosing software by the economic question}
\label{exec13:survey}

\section{Start with the banking comparison}

A bank considering a causal library faces several different purchasing and
engineering decisions. It may need a defensible effect of an invitation, a
description of plausible causal directions among funding variables, or a
policy for repeated servicing decisions. A package that produces a directed
graph does not necessarily answer the first question. A package that estimates
an average treatment effect does not necessarily answer the third. The survey
therefore takes a scientific operation as its unit of comparison: an input
with a declared observation law, a causal target, assumptions connecting the
two, and an output whose meaning survives conversion into the bank's units.

This choice changes what counts as competitive capability. Consider a request
to increase deposit funding by changing the offered rate. A system that
returns several compatible effects has provided useful information if the
observations cannot distinguish their underlying mechanisms. A system that
returns one number after choosing a convenient graph may have answered a
more restrictive question. The comparison must expose the restriction before
comparing accuracy, speed or convenience. Similarly, a symbolic expression
for an effect and an estimate of that expression are different accomplishments.
The expression identifies what information is needed; its numerical evaluation
adds statistical and computational obligations.

The proposed common request is
\begin{equation}
 q=(\mathcal U,\mathcal T,\mathcal O,\mathcal A,\mathcal Y,
       \Psi,\mathcal G,\mathcal S),
 \label{exec13:request}
\end{equation}
where $\mathcal U$ specifies the eligible population and sampling unit,
$\mathcal T$ the decision and measurement clocks, $\mathcal O$ the observed
variables and their units, and $\mathcal A$ the feasible intervention.
$\mathcal Y$ specifies the outcome and horizon. The functional $\Psi$ is the
requested causal quantity; $\mathcal G$ contains the graph or graph class and
its substantive restrictions; $\mathcal S$ describes sampling and uncertainty.
For example, one row may represent one independently sampled eligible
customer, the action may be a 25 basis point increase offered before account
opening, and the outcome may be net funding value over the next quarter.
These declarations are scientific inputs. A software interface cannot establish
them merely by accepting a table with suitable column names.

For an observational law $P$ and an admissible model class $\mathcal M$, define
\begin{equation}
 \mathcal I_\Psi(P,\mathcal M)
 =\{\Psi(M):M\in\mathcal M,\ P_M^{\rm obs}=P\}.
 \label{exec13:identified-set}
\end{equation}
This is the collection of target values compatible with the same observations
and assumptions. It is a population object. A finite-sample graph estimate,
an estimated collection of effects, a Bayesian posterior and a confidence set
each add different operations to it. A common result schema must distinguish
these objects. Reducing all four to a scalar and two bounds would make the
schema convenient at the cost of its scientific meaning.

\section{The scope of the comparison}

The landscape has two dimensions that should be searched separately. The
method dimension includes observational, latent-variable and interventional
discovery; identification and counterfactuals; heterogeneous effects;
instrumental variables and discontinuities; longitudinal treatment and survival;
policy evaluation; panels and differences in differences; synthetic controls;
time series; interference; mediation and sensitivity; transport and missingness;
Bayesian structural models; observation construction; and computational
workflow. The implementation dimension distinguishes R, Python, Java,
C/C++, other language paths and commercial offerings. A language binding,
an underlying numerical engine and a commercial wrapper can occupy different
positions in that matrix without being three independent implementations.

The publication cutoff is 7 October 2026. Acquisition time, software revision,
manual generation date and scholarly publication date are retained separately.
This matters even for apparently simple documentation: the inspected pcalg
manual identifies version 2.7-12 and a package date in 2024, while the manual
itself was generated in May 2026. The bnlearn manual identifies version 5.2.1
dated 17 July 2026. BiDAG's inspected manual was generated in May 2026 but
describes a release dated May 2023
\citep{exec13_pcalg,exec13_bnlearn,exec13_bidag}. None of these dates alone
measures current support quality, future maintenance or compatibility with
the proposed product.

A reproducible search records the query, source, date filter, returned
identities, pagination and stopping rule. For the R observational-discovery
comparison developed here, the retained bibliographic query was ``R package
causal discovery constraint score continuous DAG.'' Its saved first page
contains ten records from a reported 1,648 matches. The primary comparison
starts from the inspected pcalg and bnlearn manuals, follows the BCDAG lead,
and then follows one package-reference step to BiDAG and mcmcabn. Those five
paths define the present comparison. Ten bibliographic results are not an
exhaustive R census. The remaining methodological leads do not establish
either the presence or the absence of an R implementation.

There are two useful stopping rules. A discovery pass can stop after a
declared search budget and report incomplete coverage. A capability decision
must instead stop when the evidence is sufficient for the particular claimed
operation, or report that the decision remains unresolved. The first rule
controls search cost; the second controls the strength of the conclusion.
Conflating them would allow an arbitrary search budget to become evidence
that a missing capability does not exist. Conversely, insisting on an
unbounded census before discussing any inspected package would prevent useful
comparisons indefinitely.

Each included path needs a canonical identity, exact source or release,
license information, local source bytes, a digest, inspected technical anchors
and a claim tied to those anchors. Scholarly arguments use the original
technical text, preserved as a real paper. Software behavior uses the relevant
manual or source. Search metadata establishes identity and discovery, not
algorithmic properties. A summary prepared by the project can explain a
source but cannot substitute for an unavailable original. An unavailable
candidate remains visible with its unresolved claim and the evidence needed
to revisit it.

Conflicting evidence also requires a dated resolution. The authentic BCDAG
preprint from 2022 is useful for its Gaussian model and posterior construction.
Its contemporary package comparison cannot determine what bnlearn supports
in 2026. The inspected bnlearn manual now documents intervention and
counterfactual operations \citep{exec13_bcdag_paper,exec13_bnlearn}.
The sensible conclusion is that the comparison changed with the software.
This does not establish the validity of any particular causal use of bnlearn;
it corrects a capability description at a named version.

\section{What the five R paths actually return}

The five packages illustrate why language and algorithm labels are insufficient.
Table~\ref{exec13:r-comparison} compares their documented outputs. It is a
source-based comparison, not a benchmark of their implementations.

\begin{table}[htbp]
\centering\small
\caption{Selected R paths and the distinction their outputs require.}
\label{exec13:r-comparison}
\begin{tabularx}{\linewidth}{@{}P{.15\linewidth}XX@{}}
\toprule
Path & Documented object & Decision-relevant restriction \\
\midrule
pcalg & Partial graph or graph class; IDA effect multiset & PC discovery,
latent-variable discovery and effect calculation have distinct assumptions. \\
bnlearn & Learned network; fitted parameters; causal queries & Score search,
functional discovery and intervention queries do not share one guarantee. \\
BCDAG & Graph and parameter draws; effect posterior & Gaussian model and
graph/parameter priors define the posterior; credible intervals retain that label. \\
BiDAG & MAP graph or sampled graphs & Score, admissible search space,
parent limits and sampling configuration constrain the answer. \\
mcmcabn & Score-cache graph samples and structural queries & Arc frequency
is a structural summary; outcome intervention effects require further modeling. \\
\bottomrule
\end{tabularx}
\end{table}

In pcalg, the PC interface requires a conditional-independence test and its
significance level. Its stable skeleton option addresses ordering in the
skeleton search; it does not by itself settle all orientation conflicts.
Conservative and majority-rule variants treat ambiguous triples differently.
Some configurations can return a partial graph that is not extendable to a
DAG, and conflict-handling options can introduce bidirected marks with an
algorithmic meaning \citep{exec13_pcalg}. An importer must preserve that
meaning. Such a conflict is not automatically evidence of a latent common
cause, even though bidirected edges have a substantive interpretation in
other graphical formalisms.

A particularly consequential default is \texttt{NAdelete=TRUE}: a missing
independence-test result can cause an edge to be deleted. In the proposed
banking boundary, a failed statistical calculation should produce a diagnostic
response with its cause. It should not become evidence that two financial
variables are independent. Merely changing the option to false is also
insufficient if a failed test is then silently treated as dependence. The
wrapper must detect a missing or nonfinite result and stop that numerical
request. The package's documented flexibility and the product's selected
behavior are therefore different layers of the comparison.

The bnlearn family provides a different breadth. Its constraint-based defaults
depend on variable type, while its score-based and hybrid methods expose
different search controls. The inspected documentation gives an $\alpha$ of
0.05 for constraint procedures, locally complete observations for individual
tests with missing data, BIC defaults for common score searches, and no random
restarts by default for hill climbing. The MMHC and H2PC names denote
compositions of restriction and maximization procedures, rather than new
independent guarantees. DirectLiNGAM uses a separate functional-model route
with its own ordering and selection choices \citep{exec13_bnlearn}.
For a bank, this breadth supports useful baselines, but it also makes an
unqualified ``bnlearn result'' an inadequate description of an experiment.

BCDAG makes graph uncertainty explicit through a Gaussian DAG model and
posterior sampling. Its compatible DAG-Wishart construction is described in
the original preprint; the inspected version 1.1.4 manual distinguishes
structure-only sampling, joint graph/parameter sampling, effect summaries and
credible intervals \citep{exec13_bcdag_paper,exec13_bcdag}. The
\texttt{fast} option changes the proposal-ratio calculation through an
approximation; its default is false. A comparison must therefore specify
the graph prior, parameter hyperparameters, burn-in, retained draws, chain
configuration and approximation setting. A stable trace is useful diagnostic
evidence but is not a proof that a finite chain explored all consequential
graph regions.

BiDAG's order and partition samplers expose another tradeoff. A search space
can begin with a PC-derived skeleton and then be expanded. A MAP graph under
that space and a sample from a specified posterior serve different purposes.
Parent limits and the expansion strategy influence both computation and what
can be learned \citep{exec13_bidag}. If a true funding relationship was
excluded during search-space construction, a later apparently stable graph
sample cannot recover it without a mechanism that admits it. A benchmark
must consequently compare declared search procedures, not only the final
graphs after favorable tuning.

The mcmcabn documentation starts from a precomputed score cache, commonly
constructed with abn, and describes structural MCMC moves, priors and
structural queries \citep{exec13_mcmcabn}. This provides an alternative for
mixed node distributions and explicit structural prior beliefs. It also
introduces a scoring dependency whose assumptions and version belong to the
scientific input. The inspected documentation contains a conversion issue:
a named example places the Smoking-to-LungCancer prior at the row for the
child and column for the parent, while explanatory prose in the vignette
describes rows as parents and columns as children. A proposed integration
must resolve this discrepancy with named-node source and runtime checks.
Silently choosing one convention would be an avoidable scientific error.

These comparisons support a narrow initial choice: a pcalg route for explicitly
typed discovery and effect multisets, bnlearn as a method-specific alternative,
and Bayesian routes as separately configured sensitivity comparisons.
They do not justify porting every algorithm into a new library. The Bayesian
routes may be particularly informative when prior restrictions are part of
the scientific question; their value should be assessed by changing those
restrictions and reporting the resulting decision sensitivity.

\section{An exact funding example: identical observations, different actions}

Consider centered continuous variables $X$, $M$ and $Y$. In this constructed
example $X$ is an offered-rate margin measured in percentage points, $M$ is
a continuous funding-response index, and $Y$ is subsequent net funding value
in thousands of currency units. The index is not a binary take-up indicator.
For this first comparison, the admissible class uses the cross-sectional law
and acyclicity alone: assignment and temporal-order evidence have deliberately
been withheld. A verified banking clock would exclude any reverse model that
contradicts it. The example isolates what the observational law itself can
establish. Let the disturbances below be mutually independent standard Gaussians:
\begin{equation}
 X=\varepsilon_X,\qquad M=0.6X+\varepsilon_M,
 \qquad Y=0.8M+\varepsilon_Y.
 \label{exec13:forward}
\end{equation}
Under the causal interpretation of these equations, replacing the equation
for $X$ by a fixed assignment gives a total derivative of $0.6\times0.8=0.48$.
The observational covariance matrix, in the order $(X,M,Y)$, is
\begin{equation}
 \Sigma=\begin{pmatrix}
 1&0.6&0.48\\
 0.6&1.36&1.088\\
 0.48&1.088&1.8704
 \end{pmatrix}.
 \label{exec13:covariance}
\end{equation}
The partial covariance of $X$ and $Y$ given $M$ vanishes because
$0.48-0.6(1.088)/1.36=0$. For the Gaussian law this is the conditional
independence $X\perp Y\mid M$.

That same covariance admits the fork $X\leftarrow M\rightarrow Y$.
Take $\Var(M)=1.36$, regress $X$ on $M$ with coefficient $0.6/1.36$,
and retain the coefficient $0.8$ from $M$ to $Y$. Independent residual
variances are $1-0.6^2/1.36$ for $X$ and $1$ for $Y$. They are positive,
and the resulting covariance is exactly~\eqref{exec13:covariance}.
Intervening on $X$ now changes neither $M$ nor $Y$, so its total effect on
$Y$ is zero.

The reverse chain $X\leftarrow M\leftarrow Y$ provides a third representation.
Let $\Var(Y)=1.8704$, use coefficient $1.088/1.8704$ from $Y$ to $M$,
and $0.6/1.36$ from $M$ to $X$. The residual variance for $M$ is
$1.36-1.088^2/1.8704$, and that for $X$ is again $1-0.6^2/1.36$.
This gives the same observed law and again gives zero effect of $X$ on $Y$.
The distinct effect values across these three models are therefore
\begin{equation}
 \{0,\,0.48\}.
 \label{exec13:effect-set}
\end{equation}
The repeated zero matters if a procedure retains graph multiplicities; the
set of distinct values and a multiset must not be confused. All covariance
equalities and effects in this example have been checked with exact rational
arithmetic \citep{exec13_examples}. No discovery engine was run to obtain
them, and no finite-sample recovery rate is implied.

An observational search that preserves the undirected path $X-M-Y$ is
consistent with this ambiguity. A procedure that selects the forward chain
has supplied a graph choice in addition to the observational information.
This is precisely the distinction the IDA multiset makes useful in its
specified linear Gaussian setting \citep{exec13_pcalg}. Additional knowledge
can change the model class: a justified temporal ordering, an intervention
or a substantive structural restriction may exclude one of the alternative
graphs. Merely naming $X$ an ``offer'' does not document that restriction.
The bank must establish how the offer was assigned and when the relevant
variables were generated and observed.

The difference has an immediate monetary consequence. If $X_{\rm bp}=100X$
and $Y_{\rm currency}=1000Y$, then the forward-model derivative becomes
\begin{equation}
 \frac{\partial\E[Y_{\rm currency}\mid\doop(X_{\rm bp}=x)]}{\partial x}
 =\frac{1000}{100}\,0.48=4.8.
 \label{exec13:units}
\end{equation}
A 25 basis point change gives 120 currency units per eligible customer in
this linear construction, or 1.2 million over 10,000 such customers if the
same additive, noninterfering mechanism applies to each. The other two
compatible models give zero. These are conditional planning calculations,
not measured revenue or an investment recommendation. Aggregation and a
currency conversion cannot remove the underlying identification ambiguity.

\section{Keep graph ambiguity, statistical uncertainty and prior uncertainty apart}

The example separates uncertainties that are often compressed into one
picture. Even with the population covariance known exactly, the three causal
representations disagree about the intervention. Finite observations introduce
additional uncertainty about covariance and conditional independences.
If a graph is chosen from those observations, subsequent effect estimation
also depends on that choice. A graph-equivalence multiset is consequently
not a confidence interval with a nominal repeated-sampling coverage level.
Taking its convex hull may be useful for a conservative decision rule, but
the hull loses information about disconnected possibilities and multiplicity.

A Bayesian analysis describes another object. Write $G$ for a DAG and
$\vartheta_G$ for its parameters. The posterior distribution of an effect
can be expressed as
\begin{equation}
 \Pi(\theta\in B\mid D)
 =\sum_G \Pi(G\mid D)
   \int\ind\{\Psi(G,\vartheta_G)\in B\}
            \,d\Pi(\vartheta_G\mid G,D),
 \label{exec13:posterior}
\end{equation}
where $D$ denotes the observed data and $B$ a measurable set of effect values.
The graph prior, conditional parameter prior and likelihood jointly determine
this distribution. Compatible priors can assign the same marginal likelihood
to Markov-equivalent Gaussian DAGs, as described in the BCDAG construction
\citep{exec13_bcdag_paper}. Observational equivalence has not disappeared:
different graph weights can still reflect prior restrictions or different
admissible model classes.

The BCDAG method named \texttt{confint.bcdagCE} explicitly reports credible
intervals \citep{exec13_bcdag}. Its function name is not a reason to label
the output a frequentist confidence interval. The proposed result therefore
stores the uncertainty family, conditioning assumptions and computational
method alongside the numerical bounds. It also retains the prior and MCMC
configuration for posterior results. If a later simulation claims coverage,
that is a separate empirical statement about a specified repeated-sampling
experiment, with all failed and unavailable requests retained in its denominator.

The distinction is consequential for investment. A concentrated posterior
under a strong orientation prior can support a conditional decision while
still leaving the prior's banking justification unsettled. Conversely, an
effect collection containing zero can reveal that a proposed rate experiment
has informational value. The relevant question is then whether a feasible
experiment can distinguish the mechanisms sufficiently to change the action.
The product should make that unresolved decision visible rather than display
a more precise-looking number obtained by silently discarding alternatives.

\section{Targets must survive an identification or estimation boundary}

DoWhy provides a useful model of separating a supplied causal model,
identification and estimation. The inspected source revision is
\texttt{b06369ed} from 4 October 2026. Its graph-learning entry point is
deprecated, and its identification path operates on supplied graph,
treatment, outcome and observed-node information \citep{exec13_dowhy}.
The proposed reuse is therefore an optional identification service for an
explicit reviewed graph. It does not treat the package as an automatic
authority on the bank's diagram or as a universal discovery engine.

An exact second example shows why the target belongs to the interface.
Let $W$ be a baseline customer characteristic, $A$ a binary invitation and
$Y$ an outcome measured in thousands of currency units. Suppose
\begin{align}
 \Prob(W=1)&=\tfrac12, &
 e(W):=\Prob(A=1\mid W)&=0.2+0.6W,\\
 \E[Y(a)\mid W]&=3W+a(1+2W).&&
 \label{exec13:target-law}
\end{align}
Under consistency, exchangeability conditional on $W$, positivity and no
interference, the population average effect is
\begin{equation}
 \ATE=\sum_{w=0}^1
 \{\E[Y\mid A=1,w]-\E[Y\mid A=0,w]\}\Prob(W=w)=2.
 \label{exec13:ate}
\end{equation}
The treated group has $\Prob(W=1\mid A=1)=0.8$, while the untreated group
has $\Prob(W=1\mid A=0)=0.2$. Their average effects are consequently
$\ATT=2.6$ and $\mathrm{ATC}=1.4$. The unadjusted outcome difference is
$5-0.6=4.4$. These values are exact calculations from the declared law
\citep{exec13_examples}. The same records therefore support different
numbers for different targets, and an association can differ from all of them.

An adapter must retain whether the question concerns all eligible customers,
the invited customers or the uninvited customers. It must also retain the
contrast and the outcome scale. If $W$ is instead unobserved in the graph
$W\rightarrow A$, $W\rightarrow Y$, $A\rightarrow Y$, the stated adjustment
is unavailable. No instrument or front-door mediator has been supplied in
this example. A predictive regression on $A$ is not a repair of the missing
identification argument.

The source review suggests two additional precautions in a narrow DoWhy
boundary. Confidence intervals are not requested by default and depend on
the selected estimator. Further, the inspected identification function assigns
the argument \texttt{proceed\_when\_unidentifiable} but does not forward it
to the identifier \citep{exec13_dowhy}. This is a source observation, not
an executed package failure. The interface must validate the returned
estimand and its assumptions directly, instead of treating the presence of
that flag as a verified safeguard. An empty adjustment set, an unavailable
estimand, a model-implied zero and an engine exception need different responses.

\section{A proposed R boundary with explicit failure behavior}

The first R path is deliberately narrow: independently sampled continuous
Gaussian observations for PC, and singleton total-effect multisets for IDA
under its linear causal model. The request supplies unique variable IDs,
the actual unit of observation, sampling design, measurement clock, units,
assumptions and algorithm choices. An initial engineering design can bound
the number of variables, use one process and expose a resource budget.
Such limits are proposed workload choices, not measurements of pcalg's
maximum capability. Hitting a limit must produce an incomplete or refused
request rather than silently reducing the conditioning-set search.

Discovery may use a correlation matrix and sample size where those are
sufficient for the selected test. Effect calculation needs the original
covariance scale, matching node coordinates and a validated graph. If
$D=\operatorname{diag}(c_1,\ldots,c_p)$ changes variable units, covariance
transforms as $\Sigma_{\rm new}=D\Sigma D^{\mathsf T}$ and an effect from
coordinate $i$ to $j$ transforms by $c_j/c_i$. A correlation matrix alone
does not retain the information needed to reconstruct monetary units.
The transformation in~\eqref{exec13:units} should therefore be a routine
semantic check at the boundary.

Graph conversion needs its own type. In the documented pcalg ancestral-graph
matrix, entry $(i,j)$ describes the endpoint at $j$: zero denotes absence,
one a circle, two an arrowhead and three a tail \citep{exec13_pcalg}.
Thus $i\rightarrow j$ is represented by the pair $(2,3)$ in entries
$(i,j)$ and $(j,i)$. This rule cannot be reused indiscriminately for a binary
DAG adjacency matrix or a structural coefficient matrix. The initial PC/IDA
service does not accept a PAG merely because its entries form a square matrix.
Named endpoints, graph class and algorithmic conflicts travel separately.

A useful pre-admission oracle includes an asymmetric named chain and a column
permutation, the collider $X\rightarrow M\leftarrow Y$, the unresolved path
$X-M-Y$, and the exact covariance and unit examples above. The converter
must preserve each object's meaning. A deliberate transpose must fail the
direction check. Missing independence tests, singular inputs, invalid partial
graphs and resource exhaustion need explicit failure cases. Observations
consisting of repeated monthly rows for the same customer should receive an
unsupported-sampling response on this iid path, even when the matrix is
numerically well formed.

These are proposed adapter requirements. The exact algebraic examples have
been reproduced, but an R environment, converter and package-level oracle
run are distinct implementation evidence. Separating these claims makes it
possible to recommend a concrete prototype without treating a contract as
an installed integration. Library, command-line and MCP clients should
consume the same typed results; a display layer must not orient an unresolved
edge or relabel a missing interval merely to simplify its presentation.

\section{Reuse, porting and the next comparison}

The reuse choice should follow the selected operation and deployment need.
Using a package in an optional environment preserves upstream algorithms but
adds dependency, conversion and upgrade costs. Porting moves those obligations
into this project, including the burden of reproducing statistical behavior.
An apparent reduction in installation complexity can therefore increase
scientific maintenance. A small native reference implementation may be an
excellent oracle for a few graphs without being an adequate substitute for
a general identification or discovery library.

License and dependency evidence belong to this choice at the exact source
revision. The inspected pcalg, bnlearn and BiDAG metadata declare GPL version
2 or later; mcmcabn declares GPL version 3; BCDAG declares MIT with its
license file \citep{exec13_pcalg,exec13_bnlearn,exec13_bidag,
exec13_mcmcabn,exec13_bcdag}. DoWhy's source license permits reuse with
the specified notice, but that grant does not establish the rights or
compatibility of every dependency \citep{exec13_dowhy_license}.
Separating processes can isolate versions; it does not itself settle
distribution obligations.

The inspected DoWhy dependency declarations include Python version constraints,
NumPy greater than 2, and several numerical and causal packages. They are
constraints rather than a resolved environment, and the displayed
\texttt{0.0.0} is explicitly a dynamic-version placeholder
\citep{exec13_dowhy_project}. The proposal selects a narrow identification
path for further prototyping and defers its executable admission until a
resolved environment, dependency review and semantic oracle are available.
It does not infer performance or a support commitment from repository activity.

The same evidence structure applies to commercial comparisons. A documented
feature in a named edition can be compared with a proposed operation, but
absence from a search result cannot establish absence from the product.
A meaningful trial would hold the target, observation law, assumptions,
failed-request denominator and uncertainty requirements fixed, then compare
scientific coverage, workflow effort and total cost. Public documentation
can support a proposal before a paid trial is justified; a claim about
actual runtime behavior requires the corresponding execution evidence.

The next useful experiment is thus a matched test of a selected path, with
its alternatives and failure cases fixed in advance. For the funding example,
one question is whether the integration preserves the two distinct effect
possibilities and their units. Another is whether a specified intervention
or justified background restriction resolves a decision-relevant ambiguity.
Neither requires an enterprise deployment merely to formulate and evaluate
the prototype. Each has a clear evidential limit, an inexpensive reference
calculation and a direct connection to the eventual banking decision.
\label{exec13:survey-end}
